\documentclass[11pt]{article}
\usepackage{amsmath,amssymb,amsthm,mathrsfs}

\newtheorem{theorem}{Theorem}
\newtheorem{lemma}[theorem]{Lemma}
\newtheorem{proposition}[theorem]{Proposition}
\newtheorem{corollary}[theorem]{Corollary}

\newcommand{\N}{\mathbb N}

\title{Rigidity and finite-window failure for synchronized rounding defects}
\author{}
\date{27 August 2026}

\begin{document}
\maketitle

Let
\[
  2^{a_j}x_{j+1}=3x_j+1,
  \qquad
  2^{a_j}e_{j+1}=3e_j-1,
  \qquad a_j\in\N,
  \tag{1}
\]
where every $x_j$ is a positive integer and every $e_j$ is a real number
greater than $1$.  Put
\[
  q_j:=\frac{3^j}{2^{A_j}},
  \qquad A_j:=\sum_{i<j}a_i,
  \qquad
  m_j:=\lceil e_j\rceil,
  \qquad
  \theta_j:=m_j-e_j,
  \tag{2}
\]
and define
\[
  N_j:=x_j+m_j,
  \qquad
  d_j:=2^{a_j}m_{j+1}-3m_j+1.
  \tag{3}
\]
Then
\[
  d_j=2^{a_j}\theta_{j+1}-3\theta_j,
  \qquad -2\leq d_j\leq 2^{a_j}-1.
  \tag{4}
\]
In the normalized divergent-orbit application we also have
\[
  q_j\longrightarrow\infty,
  \qquad
  \sum_{j=0}^{\infty}\frac1{x_j}<\infty,
  \qquad
  \sum_{j=0}^{\infty}\frac1{e_j}=\infty,
  \qquad
  \sum_{\substack{j\geq0\\d_j\leq0}}\frac1{m_j}=\infty.
  \tag{5}
\]

\begin{lemma}[block defect cocycle]
Fix $s\geq0$ and $r\geq1$.  Write
\[
  B_0:=0,
  \qquad
  B_i:=\sum_{h=0}^{i-1}a_{s+h},
  \qquad
  M:=2^{B_r},
  \qquad
  P:=3^r,
  \tag{6}
\]
and put
\[
  C:=\sum_{i=0}^{r-1}3^{r-1-i}2^{B_i},
  \qquad
  R:=-\sum_{i=0}^{r-1}3^{r-1-i}2^{B_i}d_{s+i}.
  \tag{7}
\]
Then
\[
\begin{aligned}
  Mx_{s+r}&=Px_s+C,\\
  Mm_{s+r}&=Pm_s-C-R,\\
  MN_{s+r}&=PN_s-R,\\
  M\theta_{s+r}&=P\theta_s-R.
\end{aligned}
\tag{8}
\]
If $d_{s+i}\leq0$ for every $0\leq i<r$, then $R\geq0$ and
\[
  R\leq P\theta_s<R+M.
  \tag{9}
\]
Moreover, if $D:=C+M$, then
\[
  D
  =Pm_s-\bigl(R+M(m_{s+r}-1)\bigr),
  \tag{9a}
\]
and, for every $0\leq t\leq r$,
\[
\boxed{
  3^t\mid D
  \quad\Longleftrightarrow\quad
  3^t\mid R+M(m_{s+r}-1).}
  \tag{9b}
\]
The integer on the right also has the exact companion form
\[
  R+M(m_{s+r}-1)
  =P\theta_s+M(e_{s+r}-1)>0.
  \tag{9c}
\]
\end{lemma}

\begin{proof}
The first identity in (8) is the usual affine iteration of the first
recurrence in (1).  From (3) we have
\[
  2^{a_j}m_{j+1}=3m_j-1+d_j.
\]
Iterating this recurrence gives
\[
  Mm_{s+r}
  =Pm_s-C+\sum_{i=0}^{r-1}3^{r-1-i}2^{B_i}d_{s+i}
  =Pm_s-C-R.
\]
Adding the first two identities proves the identity for $N$.  Iteration of
(4) proves the identity for $\theta$.  If every displayed defect is
nonpositive, then $R\geq0$.  The inequalities $0\leq\theta_{s+r}<1$ in the
last identity of (8) give (9).

Rearranging the identity for $m$ in (8) gives (9a).  Since $P=3^r$,
reduction modulo $3^t$ proves (9b).  Finally, the last identity in (8) and
$m_{s+r}-1-\theta_{s+r}=e_{s+r}-1$ give (9c).  Its strict positivity follows
from $e_{s+r}>1$.
\end{proof}

For a nonzero integer $z$, we write $\nu_2(z):=v_2(|z|)$.

\begin{theorem}[defect-switch valuation conservation]
Assume (1)--(5), and put
\[
  c_j:=1-d_j,
  \qquad
  Z_j(c):=cx_j+m_j.
  \tag{10}
\]
There is an index $J$ such that $Z_j(c)\neq0$ for every integer $c$ and
every $j\geq J$.  For every $j\geq J$ we have
\[
\boxed{
  2^{a_j}Z_{j+1}(c_j)=3Z_j(c_j).}
  \tag{11}
\]

Partition the tail into its maximal constant-defect runs.  Let the $k$-th
run begin at $s_k$, end at $s_{k+1}$, have color $c_k$, and have total
exponent
\[
  B_k:=\sum_{j=s_k}^{s_{k+1}-1}a_j.
  \tag{12}
\]
Define
\[
\begin{aligned}
  u_k&:=\nu_2\bigl(Z_{s_k}(c_k)\bigr),\\
  v_k&:=\nu_2\bigl(Z_{s_{k+1}}(c_k)\bigr),\\
  g_k&:=\nu_2\bigl(Z_{s_{k+1}}(c_{k+1})\bigr)
       -\nu_2\bigl(Z_{s_{k+1}}(c_k)\bigr).
\end{aligned}
\tag{13}
\]
Then every run is finite,
\[
\boxed{B_k=u_k-v_k,}
\tag{14}
\]
and, for every $K\geq0$,
\[
\boxed{
  \sum_{k=0}^{K}B_k
  =u_0-v_K+\sum_{k=0}^{K-1}g_k.}
\tag{15}
\]
In particular,
\[
  \sum_{k=0}^{K-1}g_k\longrightarrow\infty.
  \tag{16}
\]

At a switch, put $c:=c_k$, $c':=c_{k+1}$,
$x:=x_{s_{k+1}}$, $m:=m_{s_{k+1}}$, and
$h:=\nu_2(c'-c)$.  If
\[
  v:=\nu_2(cx+m),
\]
then
\[
  g_k=
  \begin{cases}
    0,&v<h,\\
    h-v,&v>h,\\
    \nu_2\!\left(\dfrac{cx+m}{2^h}
      +\dfrac{c'-c}{2^h}x\right),&v=h.
  \end{cases}
  \tag{17}
\]
Consequently, a switch has positive gain if and only if
\[
  \nu_2(cx+m)=\nu_2(c'-c).
  \tag{18}
\]
\end{theorem}

\begin{proof}
Put $X_j:=x_j+1$ and $E_j:=e_j-1$.  If
$\gamma_j:=2^{a_j}-2$, then (1) gives
\[
  2^{a_j}X_{j+1}=3X_j+\gamma_j,
  \qquad
  2^{a_j}E_{j+1}=3E_j-\gamma_j.
\]
Therefore,
\[
  \frac{E_j}{X_j}-\frac{E_{j+1}}{X_{j+1}}
  =\frac{\gamma_j(E_j+X_j)}{X_j(3X_j+\gamma_j)}
  \geq0.
\]
The ratio $E_j/X_j$ has a limit.  If this limit were positive, then
$e_j\gg x_j$ on a tail, which would make $\sum_j1/e_j$ converge by (5).
Thus $E_j/X_j\to0$.  Since $x_j\to\infty$ and $m_j<e_j+1$, we obtain
\[
  \frac{m_j}{x_j}\longrightarrow0.
  \tag{19}
\]
Choose $J$ so that $m_j<x_j$ for every $j\geq J$.  If $c\geq0$, then
$Z_j(c)>0$.  If $c\leq-1$, then $Z_j(c)<0$.  Hence these linear forms do
not vanish on the chosen tail.

By the definition of $c_j$, equation (3) becomes
\[
  2^{a_j}m_{j+1}=3m_j-c_j.
\]
Multiplying the positive recurrence in (1) by $c_j$ and adding gives
\[
\begin{aligned}
  2^{a_j}Z_{j+1}(c_j)
  &=c_j(3x_j+1)+3m_j-c_j\\
  &=3Z_j(c_j),
\end{aligned}
\]
which proves (11).

Iteration through one constant-color run gives
\[
  2^{B_k}Z_{s_{k+1}}(c_k)
  =3^{s_{k+1}-s_k}Z_{s_k}(c_k).
  \tag{20}
\]
Taking $2$-adic valuations proves (14).  In particular, no constant-color
tail is possible, since its accumulated exponent would eventually exceed
the fixed valuation at its initial state.  Thus every maximal run is
finite.

By definition, $g_k=u_{k+1}-v_k$.  Summing (14) and inserting these
differences gives (15).  Since every $B_k\geq1$ and $v_K\geq0$, equation
(15) gives (16).

Finally,
\[
  c'x+m=(cx+m)+(c'-c)x.
  \tag{21}
\]
The integer $x$ is odd, so the second term on the right has valuation $h$.
If $v\neq h$, the valuation of the sum is $\min\{v,h\}$.  If $v=h$,
division by $2^h$ leaves a sum of two odd integers.  This proves (17) and
(18).
\end{proof}

\begin{corollary}[eventually constant defects are impossible]
Assume (1)--(5).  There do not exist $J\geq0$ and
$\delta\in\mathbf Z$ such that
\[
  d_j=\delta\qquad(j\geq J).
  \tag{22}
\]
In particular, none of the three nonpositive defects $-2,-1,0$ can be
eventually constant.
\end{corollary}

\begin{proof}
This is the constant-color tail excluded in the proof of the theorem.
\end{proof}

\begin{corollary}[constant-defect terminal decoder]
We use the convention $v_3(0):=\infty$.
Let $[s,s+r)$ be a constant-defect run with
\[
  d_j=1-c,
  \qquad c\in\{1,2,3\},
  \qquad s\leq j<s+r.
\]
Then
\[
\boxed{3^r\mid cx_{s+r}+m_{s+r}.}
\tag{D1}
\]
If $c\in\{1,2\}$, then
\[
\boxed{
  \min\{v_3(x_{s+r}+1),r\}
  =\min\{v_3(m_{s+r}-c),r\}.}
\tag{D2}
\]
For $c=3$, we have
\[
\boxed{
  \min\{v_3(x_{s+r}+1)+1,r\}
  =\min\{v_3(m_{s+r}-3),r\}.}
\tag{D3}
\]
In particular, a constant $d=-1$ run ending at $m_{s+r}=2$ satisfies
$3^r\mid x_{s+r}+1$.
\end{corollary}

\begin{proof}
Iteration of (11) through the run gives
\[
  2^{B_r}\bigl(cx_{s+r}+m_{s+r}\bigr)
  =3^r(cx_s+m_s).
\]
Since $2^{B_r}$ is coprime to $3$, this proves (D1).  Reducing (D1) after
writing
\[
  cx_{s+r}+m_{s+r}
  =c(x_{s+r}+1)+(m_{s+r}-c)
\]
proves (D2).  When $c=3$, the same congruence compares
$3(x_{s+r}+1)$ with $m_{s+r}-3$, which proves (D3).
\end{proof}

\begin{theorem}[switch-carry terminal decoder]
Fix a block of length $r\geq1$.  For $0\leq i\leq r$, put
\[
  x_i:=x_{s+i},
  \qquad m_i:=m_{s+i},
  \qquad c_i:=1-d_{s+i},
\]
and put
\[
  B_0:=0,
  \qquad B_i:=\sum_{h=0}^{i-1}a_{s+h},
  \qquad M:=2^{B_r}.
\]
Define the synchronized switch carry by
\[
  \mathscr S_{s,r}
  :=\sum_{i=1}^{r}
  3^{r-i}2^{B_i}(c_i-c_{i-1})x_i.
  \tag{S1}
\]
Then
\[
\boxed{
  M(c_rx_r+m_r)
  =3^r(c_0x_0+m_0)+\mathscr S_{s,r}.}
  \tag{S2}
\]
Consequently,
\[
\boxed{
  Mc_r(x_r+1)
  \equiv
  \mathscr S_{s,r}-M(m_r-c_r)
  \pmod {3^r}.}
  \tag{S3}
\]
If $c_r\neq0$, then
\[
\boxed{
  \min\{v_3(x_r+1)+v_3(c_r),r\}
  =\min\{v_3(\mathscr S_{s,r}-M(m_r-c_r)),r\}.}
  \tag{S4}
\]

Suppose that the block contains at least one color switch, and let
\[
  \ell:=\max\{1\leq i\leq r:c_i\neq c_{i-1}\}.
\]
If
\[
  3\nmid(c_\ell-c_{\ell-1})x_\ell,
  \tag{S5}
\]
then
\[
\boxed{v_3(\mathscr S_{s,r})=r-\ell.}
\tag{S6}
\]
\end{theorem}

\begin{proof}
Equation (11), followed by the change from color $c_i$ to color $c_{i+1}$,
gives
\[
\begin{aligned}
  2^{a_{s+i}}Z_{i+1}(c_{i+1})
  &=2^{a_{s+i}}Z_{i+1}(c_i)
    +2^{a_{s+i}}(c_{i+1}-c_i)x_{i+1}\\
  &=3Z_i(c_i)
    +2^{a_{s+i}}(c_{i+1}-c_i)x_{i+1}.
\end{aligned}
\]
Iteration proves (S2).  Writing
\[
  c_rx_r+m_r=c_r(x_r+1)+(m_r-c_r)
\]
and reducing (S2) modulo $3^r$ proves (S3).  Since $M$ is a unit modulo
$3^r$, comparison of truncated valuations proves (S4).

By the definition of the last switch,
\[
  \mathscr S_{s,r}
  =3^{r-\ell}
  \left(
    2^{B_\ell}(c_\ell-c_{\ell-1})x_\ell+3H
  \right)
\]
for some integer $H$.  Under (S5), the expression in parentheses is a
unit modulo $3$.  This proves (S6).
\end{proof}

\begin{corollary}[last-switch matched-shell law]
Use the notation and hypotheses of the switch-carry terminal decoder.  Assume
that $c_r\neq0$, that the block contains a color switch, and that its last
switch $\ell$ satisfies (S5).  Put
\[
  t:=r-\ell,
  \qquad
  h:=v_3(m_r-c_r),
  \qquad
  \delta:=\min\{v_3(x_r+1)+v_3(c_r),r\},
  \tag{L1}
\]
where $v_3(0):=\infty$.  If $h\neq t$, then
\[
\boxed{\delta=\min\{h,t\}.}
\tag{L2}
\]
If $h=t$, then $\delta\geq t$, and for every $1\leq q\leq r-t$ we have
\[
\boxed{
  \delta\geq t+q
  \quad\Longleftrightarrow\quad
  \frac{\mathscr S_{s,r}}{3^t}
  \equiv
  \frac{M(m_r-c_r)}{3^t}
  \pmod {3^q}.}
\tag{L3}
\]
In particular,
\[
\boxed{
  \delta>t
  \quad\Longleftrightarrow\quad
  \begin{cases}
    v_3(m_r-c_r)=t,\\[2mm]
    2^{B_\ell}(c_\ell-c_{\ell-1})x_\ell
    \equiv
    \displaystyle\frac{M(m_r-c_r)}{3^t}\pmod 3.
  \end{cases}}
\tag{L4}
\]
If $3\nmid c_r$, then $\delta=\min\{v_3(x_r+1),r\}$, so (L2)--(L4)
decode the terminal factor $x_r+1$ itself.
\end{corollary}

\begin{proof}
By (S6),
\[
  v_3(\mathscr S_{s,r})=t.
  \tag{L5}
\]
Since $M$ is a power of $2$, we also have
\[
  v_3\bigl(M(m_r-c_r)\bigr)=h.
  \tag{L6}
\]
If $h\neq t$, the two terms in the right side of (S3) have distinct
$3$-adic valuations.  Hence
\[
  v_3\bigl(\mathscr S_{s,r}-M(m_r-c_r)\bigr)=\min\{h,t\}<r,
\]
and (S4) proves (L2).

Now suppose that $h=t$.  Both terms in the same difference are divisible
by $3^t$, so (S4) gives $\delta\geq t$.  For $1\leq q\leq r-t$, the same
identity gives
\[
\begin{aligned}
  \delta\geq t+q
  &\quad\Longleftrightarrow\quad
  3^{t+q}\mid
  \mathscr S_{s,r}-M(m_r-c_r)\\
  &\quad\Longleftrightarrow\quad
  \frac{\mathscr S_{s,r}}{3^t}
  \equiv
  \frac{M(m_r-c_r)}{3^t}
  \pmod {3^q}.
\end{aligned}
\]
This proves (L3).  Finally, the factorization used in the proof of (S6)
gives
\[
  \frac{\mathscr S_{s,r}}{3^t}
  \equiv
  2^{B_\ell}(c_\ell-c_{\ell-1})x_\ell
  \pmod3.
\]
Taking $q=1$ in (L3), and combining the result with (L2), proves (L4).
The last assertion follows directly from $3\nmid c_r$.
\end{proof}

\begin{corollary}[nonpositive-block reciprocal shell]
Use the notation of the switch-carry terminal decoder, and suppose that
\[
  c_i\in\{1,2,3\}
  \qquad(0\leq i\leq r).
  \tag{N1}
\]
If the block contains a color switch, then its last switch $\ell$
automatically satisfies (S5).  Put
\[
  t:=r-\ell,
  \qquad
  \delta:=\min\{v_3(x_r+1)+v_3(c_r),r\}.
  \tag{N2}
\]
If $t\geq1$ and $\delta>t$, then
\[
\boxed{
  v_3(m_r-c_r)=t,
  \qquad
  m_r\geq3^t+c_r,
  \qquad
  \frac1{m_r}<3^{-t}.}
\tag{N3}
\]
Consequently, let a constant-color run inside the block begin at a switch
$u\geq1$ and end at $v\leq r$.  For $u\leq j\leq v$, put
\[
  t_j:=j-u,
  \qquad
  \delta_j:=\min\{v_3(x_{s+j}+1)+v_3(c_j),j\}.
\]
Then
\[
\boxed{
  \sum_{\substack{u\leq j\leq v\\t_j\geq T,\ \delta_j>t_j}}
  \frac1{m_{s+j}}
  <\sum_{t=T}^{\infty}3^{-t}
  =\frac{3^{1-T}}2
  \qquad(T\geq1).}
\tag{N4}
\]
\end{corollary}

\begin{proof}
For every $j\geq1$, the first recurrence in (1) gives
\[
  2^{a_{j-1}}x_j=3x_{j-1}+1\equiv1\pmod3.
\]
Thus $3\nmid x_j$.  At a color switch inside a block satisfying (N1), the
nonzero difference $c_i-c_{i-1}$ belongs to
$\{-2,-1,1,2\}$.  It is also a unit modulo $3$.  Since the global index
$s+\ell$ is positive, (S5) follows.

If $\delta>t$, equation (L4) gives $v_3(m_r-c_r)=t$.  When $t\geq1$, the
nonzero integer $m_r-c_r$ is a multiple of $3^t$.  Since $m_r\geq2$ and
$1\leq c_r\leq3$, it cannot be a negative multiple of $3^t$.  Hence
$m_r-c_r\geq3^t$, which proves (N3).

Apply (N3) to each block prefix ending at $s+j$.  After the switch at $u$,
the run age $t_j$ takes each nonnegative integer at most once.  Summing the
resulting bounds over the selected endpoints and then enlarging to the full
geometric tail proves (N4).
\end{proof}

\begin{corollary}[global late-shell reduction]
Assume (1)--(5).  Partition the set
\[
  \{j\geq0:d_j\leq0\}
\]
into its maximal consecutive blocks $I_b=[\alpha_b,\beta_b]$, allowing an
infinite right endpoint.  Partition each $I_b$ into maximal constant-color
runs
\[
  I_{b,k}=[u_{b,k},v_{b,k}]
  \qquad(0\leq k\leq K_b),
\]
listed in increasing order, where $K_b\in\{0,1,2,\ldots\}\cup\{\infty\}$.
For $k\geq1$ and $j\in I_{b,k}$, put
\[
  t_j:=j-u_{b,k},
  \qquad
  \delta_{b,j}:=
  \min\{v_3(x_j+1)+v_3(c_j),j-\alpha_b\}.
  \tag{G1}
\]
Choose integers $T_{b,k}\geq1$ for $k\geq1$ such that
\[
  \sum_b\sum_{k=1}^{K_b}3^{-T_{b,k}}<\infty.
  \tag{G2}
\]
Then
\[
\boxed{
  \sum_b\sum_{k=1}^{K_b}
  \sum_{\substack{j\in I_{b,k}\\
                   t_j\geq T_{b,k},\ \delta_{b,j}>t_j}}
  \frac1{m_j}<\infty.}
\tag{G3}
\]
Consequently, at least one of the following three sets has divergent
reciprocal mass:
\[
\begin{aligned}
  \mathcal E
  &:=\bigcup_b I_{b,0},\\
  \mathcal S
  &:=\bigcup_b\bigcup_{k=1}^{K_b}
    \{j\in I_{b,k}:t_j<T_{b,k}\},\\
  \mathcal L
  &:=\bigcup_b\bigcup_{k=1}^{K_b}
    \{j\in I_{b,k}:t_j\geq T_{b,k},\ \delta_{b,j}\leq t_j\}.
\end{aligned}
\tag{G4}
\]
That is,
\[
\boxed{
  \sum_{j\in\mathcal E}\frac1{m_j}=\infty,
  \quad\text{or}\quad
  \sum_{j\in\mathcal S}\frac1{m_j}=\infty,
  \quad\text{or}\quad
  \sum_{j\in\mathcal L}\frac1{m_j}=\infty.}
\tag{G5}
\]
\end{corollary}

\begin{proof}
For $k\geq1$, the run $I_{b,k}$ begins at a color switch contained in the
nonpositive block $I_b$.  Apply (N4) with $s=\alpha_b$, $u=u_{b,k}-\alpha_b$,
$v=v_{b,k}-\alpha_b$, and $T=T_{b,k}$.  We obtain
\[
  \sum_{\substack{j\in I_{b,k}\\
                   t_j\geq T_{b,k},\ \delta_{b,j}>t_j}}
  \frac1{m_j}
  <\frac{3^{1-T_{b,k}}}{2}.
\]
Summing this inequality and using (G2) proves (G3).

The four sets consisting of the summation set in (G3) and the three sets
in (G4) partition all nonpositive-defect indices.  Their total reciprocal
mass diverges by (5), whereas the first set has finite mass by (G3).
Therefore the union of the three sets in (G4) has divergent mass.  A finite
union of convergent nonnegative series is convergent, so at least one of
the three displayed series in (G5) diverges.
\end{proof}

\begin{proposition}[unbounded entrance mass in finite windows]
For every $k\geq1$, let $\lambda_k$ be the unique odd integer in
$\{1,3,5,\ldots,17\}$ such that
\[
  \lambda_k2^{k+4}\equiv13\pmod9.
  \tag{E1}
\]
Equivalently, according to the residue of $k$ modulo $6$, we have
\[
  \lambda_k=7,17,13,11,1,5.
  \tag{E2}
\]
The strict-barrier word
\[
  w_k:=(1,2,1^k)
  \tag{E3}
\]
has a positive odd integer realization and a synchronized real companion
such that
\[
  1<e_j<2,
  \qquad
  m_j=2
  \qquad(0\leq j\leq k+2),
  \tag{E4}
\]
and
\[
  (d_0,d_1,d_2,\ldots,d_{k+1})=(-1,3,-1^k).
  \tag{E5}
\]
Thus the second maximal nonpositive-defect block is one entrance run of
length $k$, and
\[
\boxed{
  \sum_{j=2}^{k+1}\frac1{m_j}=\frac{k}{2}.}
\tag{E6}
\]
In particular, entrance-run reciprocal mass has no uniform finite-window
bound, even for strict-barrier words with a positive ordinary realization.
No one nonnegative ordinary seed realizes all of the nested words $w_k$.
\end{proposition}

\begin{proof}
Since $2$ is a unit modulo $9$, equation (E1) determines one residue class
modulo $9$.  Exactly one of its representatives in $\{1,2,\ldots,17\}$ is
odd.  This proves existence and uniqueness, and direct reduction gives
(E2).

Define
\[
\begin{aligned}
  x_0&:=\frac{\lambda_k2^{k+4}-13}{9},\\
  x_1&:=\frac{\lambda_k2^{k+3}-5}{3},\\
  x_{2+i}&:=\lambda_k2^{k+1-i}3^i-1
  \qquad(0\leq i\leq k).
\end{aligned}
\tag{E7}
\]
Equation (E1) makes $x_0$ an integer.  Reducing (E1) modulo $3$ also
shows that $x_1$ is an integer.  All the integers in (E7) are positive and
odd.  Direct substitution gives
\[
  2x_1=3x_0+1,
  \qquad
  4x_2=3x_1+1,
  \qquad
  2x_{3+i}=3x_{2+i}+1
  \quad(0\leq i<k).
  \tag{E8}
\]
Since every state on the right side of a division in (E8) is odd, the
displayed powers of $2$ are exact.  Hence (E7) realizes $w_k$.

For a prefix of length $n$, the accumulated exponent is $1$ when $n=1$
and $n+1$ when $2\leq n\leq k+2$.  We have $2<3$, and
\[
  2^{n+1}<3^n
  \qquad(n\geq2),
\]
because the inequality holds at $n=2$ and its left-to-right ratio is
multiplied by $2/3$ when $n$ increases.  Thus every nonempty prefix is a
strict barrier.

We now define the real companion by
\[
\begin{aligned}
  e_{2+i}&:=1+\frac18\left(\frac23\right)^{k-i}
  &&(0\leq i\leq k),\\
  e_1&:=\frac53+\frac16\left(\frac23\right)^k,\\
  e_0&:=\frac{13}{9}+\frac19\left(\frac23\right)^k.
\end{aligned}
\tag{E9}
\]
One can check directly that
\[
  2^{a_j}e_{j+1}=3e_j-1
  \qquad(0\leq j\leq k+1).
\]
Moreover,
\[
  e_{2+i}\leq\frac98,
  \qquad
  e_1\leq\frac{11}{6},
  \qquad
  e_0\leq\frac{14}{9},
\]
and every number in (E9) is greater than $1$.  This proves (E4).

Since $m_j=2$, equation (3) becomes
\[
  d_j=2^{a_j+1}-5.
\]
The exponents in (E3) therefore give (E5), and (E6) follows.  Finally,
suppose that one nonnegative integer realized every $w_k$.  After the
first two steps, repeated use of the exponent-one recurrence would give
\[
  x_{2+i}+1=\frac{3^i}{2^i}(x_2+1)
  \qquad(i\geq0).
\]
Thus $2^i\mid x_2+1$ for every $i$, which forces $x_2=-1$.  This
contradicts nonnegativity.
\end{proof}

\begin{corollary}[run-mass valuation budget]
Let $[s,u]$ be a constant-color run with
\[
  c_j=c\in\{1,2,3\}
  \qquad(s\leq j\leq u).
\]
Put
\[
  B:=\sum_{j=s}^{u}a_j,
  \qquad
  U:=\nu_2(cx_s+m_s),
  \qquad
  V:=\nu_2(cx_{u+1}+m_{u+1}),
  \tag{V1}
\]
and define
\[
  \kappa_1:=\kappa_3:=\frac13,
  \qquad
  \kappa_2:=\frac12.
  \tag{V2}
\]
Then
\[
\boxed{
  B=U-V,
  \qquad
  \sum_{j=s}^{u}\frac1{m_j}
  \leq\kappa_cB
  =\kappa_c(U-V)
  \leq\kappa_cU.}
\tag{V3}
\]

Suppose in addition that $s\geq1$, $c_{s-1}\leq0$, and put
\[
  c^-:=c_{s-1},
  \qquad
  H:=\nu_2(c-c^-),
  \qquad
  G:=\max\{0,U-\nu_2(c^-x_s+m_s)\}.
  \tag{V4}
\]
Thus $[s,u]$ is an entrance run and $G$ is the positive part of its
switch gain.  We then have
\[
\boxed{
  U\leq H+G,
  \qquad
  H\leq a_{s-1},
  \qquad
  \sum_{j=s}^{u}\frac1{m_j}
  \leq\kappa_c(a_{s-1}+G).}
\tag{V5}
\]
Consequently, divergent reciprocal mass on entrance runs forces divergence
of the sum of the preceding exponents or of the positive entrance-switch
gains.

For the family in (E3), the second nonpositive block satisfies
\[
  c=2,
  \qquad
  B=k,
  \qquad
  U=k+2,
  \qquad
  V=2,
  \qquad
  H=2,
  \qquad
  G=k.
  \tag{V6}
\]
Hence equality holds in the reciprocal-mass estimate in (V3).
\end{corollary}

\begin{proof}
Iteration of (11) through the run gives
\[
  2^B(cx_{u+1}+m_{u+1})
  =3^{u-s+1}(cx_s+m_s).
\]
Taking $2$-adic valuations proves $B=U-V$.  Since every $a_j\geq1$, the
run length is at most $B$.  If $c\in\{1,3\}$, parity in
\[
  2^{a_j}m_{j+1}=3m_j-c
\]
shows that $m_j$ is odd.  Since $m_j\geq2$, we have $m_j\geq3$.  If
$c=2$, the same identity shows that $m_j$ is even, so $m_j\geq2$.
Summing these pointwise bounds proves (V3).

At the entrance switch,
\[
  cx_s+m_s=(c^-x_s+m_s)+(c-c^-)x_s.
  \tag{V7}
\]
Since $x_s$ is odd, the second term on the right has valuation $H$.  If
the two terms in (V7) have distinct valuations, then $U$ is their minimum
and hence $U\leq H$.  If they have the same valuation, then this common
valuation is $H$ and $G=U-H$.  In both cases, $U\leq H+G$.

Write $d_{s-1}=1-c^-\geq1$.  Equation (4) gives
\[
  d_{s-1}\leq2^{a_{s-1}}-1.
\]
Therefore,
\[
  1\leq c-c^-=c+d_{s-1}-1\leq2^{a_{s-1}}+1<2^{a_{s-1}+1}.
\]
It follows that $H\leq a_{s-1}$.  Combining this with (V3) proves (V5).
After summing (V5), a divergent left side forces at least one of the two
nonnegative sums on the right to diverge.

Finally, equations (E7) and (E4) give
\[
  \nu_2(2x_2+m_2)=k+2,
  \qquad
  \nu_2(2x_{k+2}+m_{k+2})=2.
\]
At the entrance, $c^-=-2$, so (V4) gives $H=2$ and $G=k$.  The remaining
claims in (V6) follow from (E3), and (E6) gives equality in (V3).
\end{proof}

\begin{corollary}[negative shadows of defect runs]
Let $[s,u]$ be a constant-color run with
\[
  c_j=c\in\{1,2,3\}
  \qquad(s\leq j\leq u).
  \tag{Y1}
\]
The following nonempty sequences consist of positive odd integers and are
accelerated $3y-1$ paths:
\[
\begin{array}{c|c|c}
  c&\text{shadow states}&\text{discarded ceiling index}\\ \hline
  1&y_j:=m_j\quad(s\leq j\leq u)&\varnothing\\
  2&y_j:=m_j/2\quad(s\leq j<u)&u\\
  3&y_j:=m_j/3\quad(s<j\leq u)&s.
\end{array}
\tag{Y2}
\]
Thus every consecutive pair of displayed shadow states satisfies
\[
\boxed{
  2^{a_j}y_{j+1}=3y_j-1,
  \qquad
  a_j=v_2(3y_j-1).}
\tag{Y3}
\]

Let $H:=10.789$, the certified reciprocal-mass bound for every finite
positive simple accelerated $3y-1$ path.  If the corresponding nonempty
shadow in (Y2) is simple, then
\[
\boxed{
  \sum_{j=s}^{u}\frac1{m_j}
  <
  \begin{cases}
    H,&c=1,\\[1mm]
    H/2+1/2,&c=2,\\[1mm]
    H/3+1/3,&c=3.
  \end{cases}}
\tag{Y4}
\]
The same bounds hold when the indicated shadow is empty.  Consequently,
if a run violates its corresponding bound in (Y4), its shadow contains a
repeated state and hence a positive cycle of the accelerated $3y-1$ map.
\end{corollary}

\begin{proof}
On the run, the ceiling recurrence is
\[
  2^{a_j}m_{j+1}=3m_j-c.
  \tag{Y5}
\]
If $c=1$, parity in (Y5) shows that every $m_j$ in the run is odd, and
(Y3) follows directly with $y_j=m_j$.

Suppose that $c=2$.  Parity in (Y5) shows that $m_j$ is even at every
index of the run.  For $j<u$, put $m_j=2y_j$.  Whenever $j+1<u$, division
of (Y5) by $2$ gives
\[
  2^{a_j}y_{j+1}=3y_j-1.
\]
The same identity with the next ceiling state still even shows that each
displayed $y_j$ is odd.  Hence the powers $2^{a_j}$ are exact.

Finally, suppose that $c=3$.  Parity in (Y5) shows that every $m_j$ in the
run is odd.  For $s<j\leq u$, the preceding instance of (Y5) gives
\[
  2^{a_{j-1}}m_j=3(m_{j-1}-1),
\]
so $3\mid m_j$.  Therefore $y_j=m_j/3$ is a positive odd integer, and
division of (Y5) by $3$ proves (Y3).  This establishes (Y2)--(Y3).

If $c=1$, the identity $m_j=y_j$ and the certified simple-path bound give
the first line of (Y4).  If $c=2$, the shadow contributes less than $H/2$
to the ceiling reciprocal mass, while the discarded endpoint satisfies
$m_u\geq2$.  If $c=3$, the shadow contributes less than $H/3$, while
$m_s$ is an odd integer at least $3$.  These observations prove (Y4),
including the empty-shadow cases.

If a run violates (Y4), the corresponding shadow cannot be simple.  It
therefore contains two equal positive odd states.  The accelerated map
\[
  y\longmapsto\frac{3y-1}{2^{v_2(3y-1)}}
\]
is deterministic, so the segment between those equal states is a positive
cycle.
\end{proof}

\begin{proposition}[bounded aperiodic all-nonpositive countermodel]
There is an aperiodic infinite strict-barrier valuation word with the
following properties.

\begin{enumerate}
  \item Its coefficient $q_j$ tends to infinity.
  \item It has one fixed real companion satisfying
  \[
    4<e_j\leq14,
    \qquad d_j\in\{-1,0\}\qquad(j\geq0).
    \tag{C1}
  \]
  In particular, $\sum_j1/e_j=\infty$.
  \item Every finite prefix has a positive ordinary integer realization.
  \item The least positive dyadic cylinder representatives of its prefixes
  tend to infinity.  Hence the complete word is not realized by one
  nonnegative ordinary integer.
\end{enumerate}
\end{proposition}

\begin{proof}
We use the following three length-$16$ words:
\[
\begin{aligned}
  w_U&:=1113(12)^6,\\
  w_V&:=11221122(12)^4,\\
  w_W&:=1122121122(12)^3.
\end{aligned}
\tag{C2}
\]
Every word in (C2) has total exponent $24$.  Their exponent-prefix vectors
are
\[
\begin{aligned}
 &(1,2,3,6,7,9,10,12,13,15,16,18,19,21,22,24),\\
 &(1,2,4,6,7,8,10,12,13,15,16,18,19,21,22,24),\\
 &(1,2,4,6,7,9,10,11,13,15,16,18,19,21,22,24).
\end{aligned}
\tag{C3}
\]
Coordinatewise, these vectors are at most
\[
  (1,3,4,6,7,9,11,12,14,15,17,19,20,22,23,25)
  =\bigl(\lfloor k\log_2 3\rfloor\bigr)_{k=1}^{16}.
  \tag{C4}
\]
Thus every nonempty prefix of each word is a strict barrier.  Also
\[
  P:=3^{16}=43046721>16777216=2^{24}=:M.
  \tag{C5}
\]

Each word returns the companion ceiling to $m=5$.  The complete ceiling
state sequences are
\[
\begin{aligned}
  U:&\quad(5,7,10,14,5,(7,5)^6),\\
  V:&\quad(5,7,10,7,5,7,10,7,5,(7,5)^4),\\
  W:&\quad(5,7,10,7,5,7,5,7,10,7,5,(7,5)^3).
\end{aligned}
\tag{C6}
\]
The corresponding defect words are
\[
\begin{aligned}
  \delta_U&:=(0,0,-1,-1,0^{12}),\\
  \delta_V&:=(0,0,-1,0,0,0,-1,0^9),\\
  \delta_W&:=(0,0,-1,0^5,-1,0^7).
\end{aligned}
\tag{C7}
\]
Indeed, direct substitution in
\[
  2^{a_j}m_{j+1}=3m_j-1+d_j
\]
verifies (C6) and (C7).

Let $\theta$ be the ceiling error at a block boundary.  Iterating
$2^{a_j}\theta_{j+1}=3\theta_j+d_j$ gives a forward affine map on each
block.  The admissible starting interval, the resulting image interval,
and the inverse map are as follows:
\[
\begin{array}{c|c|c|c}
 &\text{domain}&\text{image}&\phi_i(t)\\ \hline
 U&[20/81,4/9)&[0,531441/1048576)&(Mt+10628820)/P\\
 V&[580/2187,4/9)&[0,964467/2097152)&(Mt+11416140)/P\\
 W&[4964/19683,4/9)&[0,1034451/2097152)&(Mt+10856268)/P.
\end{array}
\tag{C8}
\]
For completeness, these intervals follow from the prefix identity
\[
  2^{B_i}\theta_i=3^i\theta_0-R_i.
  \tag{C9}
\]
At every prefix we impose $0\leq\theta_i<1$, take the maximum of the lower
bounds $R_i/3^i$, and take the minimum of the upper bounds
$(R_i+2^{B_i})/3^i$.  This gives (C8) by rational arithmetic.

Put
\[
\begin{aligned}
  p_U&:=\frac{2125764}{5253901},\\
  p_V&:=\frac{78732}{181169},\\
  p_W&:=\frac{638604}{1545265},
\end{aligned}
\qquad
  I:=[p_U,p_V].
\tag{C10}
\]
These are the respective fixed points of the three inverse maps in (C8),
and
\[
  p_U<p_W<p_V.
\]
Direct cross multiplication in (C8) shows that $I$ lies in all three
domains and all three image intervals.  Each $\phi_i$ is increasing, has
contraction ratio $M/P<1$, and has its fixed point in $I$.  Therefore
\[
  \phi_i(I)\subseteq I
  \qquad(i\in\{U,V,W\}).
  \tag{C11}
\]

Let $\omega=(\omega_0,\omega_1,\ldots)$ be any infinite word over
$\{U,V,W\}$.  For every $k$, the limit
\[
  \theta_k
  :=\lim_{L\to\infty}
  \phi_{\omega_k}\circ\phi_{\omega_{k+1}}\circ\cdots
  \circ\phi_{\omega_L}(p_U)
  \tag{C12}
\]
exists in $I$, since all three maps have the same contraction ratio less
than $1$.  It satisfies
\[
  \theta_k=\phi_{\omega_k}(\theta_{k+1}).
\]
Thus concatenating the blocks selected by $\omega$ gives one fixed real
companion with the ceiling states and defects in (C6) and (C7).  This proves
(C1).  Every block is a strict barrier and multiplies the coefficient by
$P/M>1$, so the complete concatenation is a strict barrier and $q_j\to
\infty$.

It remains to choose the driver so that it is aperiodic and has no ordinary
integer seed.  Enumerate all eventually periodic ternary words as
$\eta^{(0)},\eta^{(1)},\ldots$.  Suppose a finite driver has been chosen,
with valuation prefix $v$ and least positive cylinder representative
$n(v)$.  The three extensions $vw_U,vw_V,vw_W$ have the same length and
total exponent, and their exact cylinders are pairwise disjoint.  Hence
$n(v)$ realizes at most one of them.  At least two choices have a new least
representative different from $n(v)$.  At stage $k$, choose one of those
two symbols which is also different from the $k$-th symbol of
$\eta^{(k)}$.

The resulting driver differs from every eventually periodic ternary word,
so it is aperiodic.  Since the three codewords are distinct and have the
same length, eventual periodicity of their concatenation would force the
aligned codeword sequence to be eventually periodic.  Hence the valuation
word is also aperiodic.  Cylinder compatibility gives
\[
  n(vw_i)=n(v)+h\,2^{A(v)+1}
  \qquad(h\geq1)
  \tag{C13}
\]
at every stage.  Therefore the least representatives tend to infinity.
Every finite prefix is nevertheless realized by its own least positive
representative.  Finally, $e_j\leq14$ gives
$\sum_j1/e_j=\infty$.
\end{proof}

\begin{proposition}[arbitrarily long zero-defect finite windows]
For every $k\geq1$, the strict-barrier word
\[
  w_k:=(1,2)^k
  \tag{23}
\]
is realized by a positive integer orbit segment and by a synchronized real
companion for which
\[
  e_{2i}=5,
  \qquad e_{2i+1}=7,
  \qquad d_j=0
  \quad(0\leq j<2k).
  \tag{24}
\]
Nevertheless, if $D(w):=C(w)+2^{A(w)}$, then
\[
  v_3(D(w_k))=0.
  \tag{25}
\]
Thus the terminal $3$-adic deficit of this block is its full length $2k$.
\end{proposition}

\begin{proof}
For $0\leq i\leq k$, define
\[
  x_{2i}:=2\cdot9^i8^{k-i}-5,
  \tag{26}
\]
and, for $0\leq i<k$, define
\[
  x_{2i+1}:=3\cdot9^i8^{k-i}-7.
  \tag{27}
\]
Every displayed integer is positive.  If
$N_i:=2\cdot9^i8^{k-i}$, then
\[
  3x_{2i}+1=3N_i-14\equiv2\pmod4,
\]
so the valuation at $x_{2i}$ is exactly $1$.  Moreover
\[
  3x_{2i+1}+1=9^{i+1}8^{k-i}-20\equiv4\pmod8,
\]
so the valuation at $x_{2i+1}$ is exactly $2$.  Division gives (26) and
(27) at the following states.  Hence the segment realizes $w_k$ exactly.

The companion values in (24) satisfy
\[
  \frac{3\cdot5-1}{2}=7,
  \qquad
  \frac{3\cdot7-1}{4}=5.
\]
They are integers, so $m_j=e_j$, $\theta_j=0$, and (3) gives $d_j=0$.
Adding the two paired recurrences also gives
\[
  x_j+e_j=q_j(x_0+5)
\]
throughout the segment.

For every even prefix of length $2i$ we have
$2^{A}=8^i<9^i=3^{2i}$.  For every odd prefix of length $2i+1$ we have
\[
  2^{A}=2\cdot8^i<3\cdot9^i=3^{2i+1}.
\]
Thus every prefix is a strict barrier.

For the block $(1,2)$, the coefficient pair and carry are
\[
  P=9,
  \qquad M=8,
  \qquad C=5.
\]
The concatenation formula therefore gives
\[
  C(w_k)=5\frac{9^k-8^k}{9-8}=5(9^k-8^k).
\]
It follows that
\[
  D(w_k)=5\cdot9^k-4\cdot8^k
  \equiv-4(-1)^k\not\equiv0\pmod3.
\]
This proves (25).
\end{proof}

\begin{corollary}[no finite-window bridge to deficit two or three]
For every $L\geq1$, there is a paired strict-barrier block of length at
least $L$ whose rounding defects are all nonpositive, but whose terminal
$3$-adic deficit is greater than $3$.  Therefore the hypotheses of a valid
bridge from nonpositive rounding-defect mass to a deficit-two or
deficit-three terminal congruence must use one fixed ordinary orbit at
unboundedly many depths.  No implication from the corresponding properties
of one finite block can supply that bridge.
\end{corollary}

\begin{proof}
Choose $k\geq2$ with $2k\geq L$ and apply the proposition.  Its terminal
deficit is $2k\geq4$.
\end{proof}

\end{document}
